\documentclass[11pt]{article}
\pagestyle{plain}

% \pgfplotsset{compat=1.18}
\usepackage[margin=1in]{geometry}
\usepackage{amsmath,amsfonts,amssymb,amsthm}
\usepackage[ruled,vlined,linesnumbered,commentsnumbered]{algorithm2e}
\SetKwInput{KwIn}{Input}
\SetKwInput{KwOut}{Output}
\SetKwInput{KwGlobal}{Global variable}
\SetKwInput{KwOracle}{Oracle}
\SetKwComment{Comment}{// }{}
\usepackage{graphicx}
\usepackage{enumerate}
\usepackage{bbm}
\usepackage{verbatim}
\usepackage{hyperref,color}
\usepackage[capitalize,nameinlink]{cleveref}
\usepackage[dvipsnames]{xcolor}
\hypersetup{
	colorlinks=true,
	pdfpagemode=UseNone,
	citecolor=OliveGreen,
	linkcolor=NavyBlue,
	urlcolor=Magenta,
	pdfstartview=FitW
}
\usepackage{appendix}
\usepackage{makecell}
\usepackage{tablefootnote}
\crefname{appsec}{Appendix}{Appendices}
\usepackage{tikz}
\usepackage{pgfplots}
\usepackage{xifthen}
\usepackage{subcaption}
\usepackage{hhline}
\usepackage{listings}

\theoremstyle{plain}
\newtheorem{theorem}{Theorem}[section]
\newtheorem{proposition}[theorem]{Proposition}
\newtheorem{lemma}[theorem]{Lemma}
\newtheorem{corollary}[theorem]{Corollary}
\newtheorem{conjecture}[theorem]{Conjecture}
\newtheorem{problem}[theorem]{Problem}
\newtheorem{claim}[theorem]{Claim}
\newtheorem{fact}[theorem]{Fact}

\theoremstyle{definition}
\newtheorem{definition}[theorem]{Definition}
\newtheorem{example}[theorem]{Example}
\newtheorem{observation}[theorem]{Observation}
\newtheorem{assumption}[theorem]{Assumption}
\newtheorem*{assumption*}{Assumption}
\newtheorem{condition}[theorem]{Condition}

\theoremstyle{remark}
\newtheorem{remark}[theorem]{Remark}

\crefname{lemma}{Lemma}{Lemmas}
\crefname{theorem}{Theorem}{Theorems}
\crefname{definition}{Definition}{Definitions}
\crefname{fact}{Fact}{Facts}
\crefname{claim}{Claim}{Claims}
\crefname{proposition}{Proposition}{Propositions}


\newcommand{\dif}{\,\mathrm{d}}
%\newcommand{\E}{\mathbb{E}}
%\newcommand{\Var}{\mathrm{Var}}
\newcommand{\Cov}{\mathrm{Cov}}
\newcommand{\MI}{\mathrm{I}}
%\newcommand{\Ent}{\mathrm{Ent}}
\newcommand{\ind}{\mathbbm{1}}
\newcommand{\allone}{\mathbf{1}}
\newcommand{\tr}{\mathrm{tr}}
\newcommand{\wt}{\mathrm{weight}}
\DeclareMathOperator*{\argmax}{arg\,max}
\newcommand{\norm}[1]{\left\lVert #1 \right\rVert}
\newcommand{\TV}[2]{\left\|#1 - #2\right\|_{\mathrm{TV}}}
\newcommand{\ceil}[1]{\left\lceil #1 \right\rceil}
\newcommand{\floor}[1]{\left\lfloor #1 \right\rfloor}
\newcommand{\ip}[2]{\left\langle #1 , #2 \right\rangle}
\newcommand{\grad}{\nabla}
\newcommand{\hessian}{\nabla^2}
\newcommand{\trans}{\intercal}
\newcommand{\diag}{\mathrm{diag}}
\newcommand{\poly}{\mathrm{poly}}
\newcommand{\dist}{\mathrm{dist}}
\newcommand{\sgn}{\mathrm{sgn}}
\newcommand{\fpras}{\mathsf{FPRAS}}
\newcommand{\fpaus}{\mathsf{FPAUS}}
\newcommand{\fptas}{\mathsf{FPTAS}}

\newcommand{\eps}{\varepsilon}

\newcommand{\N}{\mathbb{N}}
\newcommand{\Z}{\mathbb{Z}}
\newcommand{\Q}{\mathbb{Q}}
\newcommand{\R}{\mathbb{R}}
\newcommand{\C}{\mathbb{C}}

\renewcommand{\AA}{\mathcal{A}}
\newcommand{\BB}{\mathcal{B}}
\newcommand{\CC}{\mathcal{C}}
\newcommand{\DD}{\mathcal{D}}
\newcommand{\EE}{\mathcal{E}}
\newcommand{\FF}{\mathcal{F}}
\newcommand{\GG}{\mathcal{G}}
\newcommand{\HH}{\mathcal{H}}
\newcommand{\II}{\mathcal{I}}
\newcommand{\JJ}{\mathcal{J}}
\newcommand{\KK}{\mathcal{K}}
\newcommand{\LL}{\mathcal{L}}
\newcommand{\MM}{\mathcal{M}}
\newcommand{\NN}{\mathcal{N}}
\newcommand{\OO}{\mathcal{O}}
\newcommand{\PP}{\mathcal{P}}
\newcommand{\QQ}{\mathcal{Q}}
\newcommand{\RR}{\mathcal{R}}
\renewcommand{\SS}{\mathcal{S}}
\newcommand{\TT}{\mathcal{T}}
\newcommand{\UU}{\mathcal{U}}
\newcommand{\VV}{\mathcal{V}}
\newcommand{\WW}{\mathcal{W}}
\newcommand{\XX}{\mathcal{X}}
\newcommand{\YY}{\mathcal{Y}}
\newcommand{\ZZ}{\mathcal{Z}}

\newcommand{\KL}[2]{D_{\mathrm{KL}} \left( #1 \,\Vert\, #2 \right)}
\newcommand{\KLbig}[2]{D_{\mathrm{KL}} \left( #1 \,\big\Vert\, #2 \right)}
\newcommand{\KLBig}[2]{D_{\mathrm{KL}} \left( #1 \,\Big\Vert\, #2 \right)}

\newcommand{\SKL}[2]{D_{\mathrm{SKL}} \left( #1 \,\Vert\, #2 \right)}
\newcommand{\SKLbig}[2]{D_{\mathrm{SKL}} \left( #1 \,\big\Vert\, #2 \right)}
\newcommand{\SKLBig}[2]{D_{\mathrm{SKL}} \left( #1 \,\Big\Vert\, #2 \right)}

\newcommand{\qandq}{\quad\text{and}\quad}
\newcommand{\supp}{\mathrm{supp}}

\newcommand{\DKL}[2]{\-D_{\-{KL}}\left(#1 \parallel #2\right)}
\newcommand{\DTV}[2]{\-D_{\mathrm{TV}}\left({#1},{#2}\right)}
% \newcommand{\dist}{\mathrm{dist}}
\newcommand{\e}{\mathrm{e}}
\renewcommand{\epsilon}{\varepsilon}
\renewcommand{\emptyset}{\varnothing}
% \newcommand{\norm}[1]{\left\Vert#1\right\Vert}
\newcommand{\set}[1]{\left\{#1\right\}}
\newcommand{\tuple}[1]{\left(#1\right)} 
% \newcommand{\eps}{\varepsilon}
\newcommand{\inner}[2]{\left\langle #1,#2\right\rangle}
\newcommand{\tp}{\tuple}
\newcommand{\ol}{\overline}
\newcommand{\abs}[1]{\left\vert#1\right\vert}
\newcommand{\ctp}[1]{\left\lceil#1\right\rceil}
\newcommand{\ftp}[1]{\left\lfloor#1\right\rfloor}
\newcommand{\Ind}[1]{\-{Ind}_{#1}}
\newcommand{\todo}[1]{{\color{red} TODO: #1}}

\newcommand{\Cor}{\Psi^{\mathrm{sym}}}
\newcommand{\cor}{\Psi^{\-{cor}}}
\newcommand{\Inf}{\Psi}
\newcommand{\MAP}{\-{MAP}}

\def\*#1{\boldsymbol{#1}} % Use \*A for \mathbf{A}
\def\+#1{\mathcal{#1}} % Use \+A for \mathcal{A}
\def\-#1{\mathrm{#1}} % Use \-A for \mathrm{A}
\def\=#1{\mathbb{#1}} % Use \^A for \mathbb{A}
\def\!#1{\mathfrak{#1}} % Use \!A for \mathfrak{A}

\newcommand*\diff{\mathop{}\!\mathrm{d}}

% probability
% \DeclareMathOperator*{\oPr}{\mathbf{Pr}}
\def\oPr{\mathop{\mathrm{Pr}}}
\renewcommand{\Pr}[2][]{ \ifthenelse{\isempty{#1}}
  {\oPr\left[#2\right]}
  {\oPr_{#1}\left[#2\right]} } % Use \Pr[a]{b} for \mathbf{Pr}_a[b], \Pr{b} for  \mathbf{Pr}[b]

% expectation: relies on the package "dsfont"
% \DeclareMathOperator*{\oE}{\mathbb{E}}
\def\oE{\mathop{\mathbb{E}}}
\newcommand{\E}[2][]{ \ifthenelse{\isempty{#1}}
  {\oE\left[#2\right]}
  {\oE_{#1}\left[#2\right]} }

% variance
%\DeclareMathOperator*{\oVar}{\mathbf{Var}}
\def\oVar{\mathrm{Var}}
\newcommand{\Var}[2][]{ \ifthenelse{\isempty{#1}}
  {\oVar\left[#2\right]}
  {\oVar_{#1}\left[#2\right]} }

% entropy
% \DeclareMathOperator*{\oEnt}{\mathbf{Ent}}
\def\oEnt{\mathrm{Ent}}
\newcommand{\Ent}[2][]{ \ifthenelse{\isempty{#1}}
  {\oEnt\left[#2\right]}
  {\oEnt_{#1}\left[#2\right]} }

\newcommand{\PhiEnt}[2][]{ \ifthenelse{\isempty{#1}}
  {\oEnt^\phi\left[#2\right]}
  {\oEnt^\phi_{#1}\left[#2\right]} }


\newenvironment{Exampleparagraph}[1][]
{
    \par 
    \noindent 
    \textbf{\color{blue} #1}
    \par 
    \vspace{0.5em} 
    \itshape 
}
{
    \par 
    \vspace{1em} 
}

%\input{local-macro}

\newcommand{\RED}[1]{\textcolor{red}{#1}}

\newcommand{\bern}[1]{\mathrm{Bern}\left(#1\right)}
\newcommand{\pll}{\parallel}



\title{CS 7545 F26, Homework 3}
\date{Due on November 12, 2026}
\author{{\textbf{Zejia Chen}} \\ \textbf{GTID: 904178159}}

\pgfplotsset{compat=1.18} 
\begin{document}

\maketitle

\section*{Exercise 1}
\textbf{Part A:}
Since $\E[z]{z}=0, \E[z]{z^2}=1$, we have the following equalities
$$
  L(\theta) = \E[z]{(z-\theta)^2} + \lambda(\theta^2 - 1)^2
  = \E[z]{z^2 - 2z\theta} + \theta^2 + \lambda(\theta^2 - 1)^2
  = 1 + \theta^2 + \lambda(\theta^2 - 1)^2 ;
$$
\begin{align*}
  \hat L_S(\theta) &= \frac 1n \sum_{i=1}^n \tp{z_i^2-2z_i\theta + \theta^2 + \lambda(\theta^2-1)^2}
  = \tp{\frac 1n\sum_{i=1}^n z_i^2}- 2\overline{z}_S\theta + \theta^2 + \lambda(\theta^2-1)^2
  \\&=\tp{\frac 1n\sum_{i=1}^n z_i^2} + (\theta-\overline z_S)^2 - \overline{z}_S^2 + \lambda(\theta^2-1)^2
\end{align*}
\textbf{Part B:}
\\1. With $\lambda = \frac 12$, $L(\theta) = \frac{1}{2}(\theta^4+3)$ and $L''(\theta) = 6\theta^2$.
Therefore, the minimizer $\theta^\star = 0$ and $L''(\theta^\star) = 0$.
\\2. $\hat L_S'(\theta) = 2(\theta^3 - \overline z_S)$.
Since $\hat L_S'(\theta)$ is strictly increasing, the ERM is unique and $\hat L_S'(\hat \theta_S) = 0$ implies that  $\hat \theta_S = (\overline z_S)^{1/3}$.
\\3. We have that $L(\theta^\star) = 3/2$ and 
$$
  \E{L(\hat \theta_S) - L(\theta^\star)} = \E{1 + (\overline z_S)^{2/3} + \frac 12\tp{(\overline z_S)^{2/3}-1}^2} - \frac 32
  = \frac 12\E{(\overline z_S)^{4/3}} = \frac 12\E{\abs{\overline z_S}^{4/3}} .
$$
Since $x\mapsto x^{4/3}$ is convex, by Jensen's inequality, $\E{L(\hat \theta_S) - L(\theta^\star)} = \frac 12\E{\abs{\overline z_S}^{4/3}} \geq \frac 12\E{\abs{\overline z_S}}^{4/3}$.
\\4. Since $z_i\sim \+N(0,1)$, $\overline z_S \sim \+N(0,1/n)$. Therefore, the expected excess loss satisfies that 
$$
  \E{L(\hat \theta_S) - L(\theta^\star)} \geq \frac 12\E{\abs{\overline z_S}}^{4/3} = \frac 12\tp{\sqrt{\frac 2{\pi n}}}^{4/3} = \Omega\tp{n^{-2/3}} .
$$
\textbf{Part C:} 
1. $L(\theta) = \frac{2}{3}(\theta^2 - \frac 14)^2 + \frac{13}8$.
Therefore, the global minimizers satisfy $(\theta^\star)^2 = \frac 14$, and hence, $\theta^\star = \pm \frac 12$.
\\2. Direct calculation shows that $\hat L_S(\theta) - \hat L_S(-\theta) = (\theta - \overline z_S)^2 -  (\theta + \overline z_S)^2 = -4\overline z_S\theta$.
Since $\hat \theta_S$ is the minimizer, $\hat L_S(-\hat \theta_S) - \hat L_S(\hat \theta_S) = 4\overline{z}_S\hat\theta_S \geq 0$.
Also, $\hat L_S'(0) = -2\overline z_S \neq 0$, implying $\hat \theta_S\neq 0$.
Therefore, $\hat\theta_S$ has the same sign as $\overline z_S$.
\\3. $\hat L'_S(\theta) = \frac 83 \theta^3 - \frac 23\theta - 2\overline z_S$.
Since $\hat \theta_S$ is the minimizer, $\hat L'_S(\hat \theta_S)=0$, that is, $4\hat \theta_S^3 - \hat \theta_S = 3\overline z_S$.
If $\overline z_S$ is positive, then $\sgn(\hat \theta_S) > 0$ and $\hat \theta_S \geq \frac 12$.
In this case, $3\overline z_S = (2\hat \theta_S+1)(2\hat \theta_S-1)\hat \theta_S \geq 2(\hat \theta_S-\frac 12)$.
If $\overline z_S$ is negative, then $\sgn(\hat \theta_S) < 0$ and $\hat \theta_S \leq -\frac 12$.
In this case, $3\overline z_S = (2\hat \theta_S+1)(2\hat \theta_S-1)\hat \theta_S \leq 2(\hat \theta_S+\frac 12)$.
Hence, $\abs{\hat \theta_S - \frac 12 \sgn(\overline z_S)}\leq \frac 32 \abs{\overline z_S}$.
\\4. By Gaussian tail bound, $\Pr{\abs{\overline z_S} \geq t} \leq 2\exp(-nt^2/2)$. By $\abs{\hat \theta_S - \frac 12 \sgn(\overline z_S)}\leq \frac 32 \abs{\overline z_S}$, we conclude that
$$
  \Pr{\abs{\hat \theta_S - \frac 12 \sgn(\overline z_S)}\leq \frac 32 t} \geq \Pr{\abs{\overline z_S}\leq t} \geq 1-2\exp\tp{-nt^2/2}.
$$
By taking $t = \sqrt{2\ln(2/\delta)/n}$, we complete the proof.
\section*{Exercise 2}
We have that $d(x,x) = \sqrt{2k(x,x)-2k(x,x)} = 0$.
The symmetry of $d$ follows from the symmetry of $k$.
Define the matrix $K \in \R^{\+X\times \+X}$ and $K = (k(x,y))_{x,y\in \+X}$ is PSD.
Define $\*1_{a,b}$ to be the vector satisfying that $\*1_{a,b}(x) = \ind[x=a] - \ind[x=b]$. 
Then $d(x,y) = \sqrt{\*1_{x,y}^\top K\*1_{x,y}}$.
For the triangle inequality, we have that
\begin{align*}
  d(x,y) &= \sqrt{\*1_{x,y}^\top K\*1_{x,y}}
    = \sqrt{(\*1_{x,z}+\*1_{z,y})^\top K(\*1_{x,z}+\*1_{z,y})}
    \\&= \sqrt{\*1_{x,z}^\top K \*1_{x,z} + \*1_{z,y}^\top K \*1_{z,y} + 2\inner{K^{1/2}\*1_{x,z}}{K^{1/2}\*1_{z,y}}}
    \\&\leq \sqrt{\norm{K^{1/2}\*1_{x,z}}_2^2 + \norm{K^{1/2}\*1_{z,y}}_2^2 + 2\norm{K^{1/2}\*1_{x,z}}_2\norm{K^{1/2}\*1_{z,y}}_2}
    \\&=d(x,z)+d(z,y),
\end{align*}
where the inequality follows from the Cauchy-Schwarz inequality.
\section*{Exercise 3}
Let $K = (k(x_i,x_j))_{i,j\in[n]}$. Then by Section 3.4 in Lecture notes 9, the prediction on a test point $x_0$ is 
\begin{equation}\label{eq:E3}
  k_{x_0}^\top (K+n\lambda I)^{-1}y\qquad \text{where}\qquad (k_{x_0})_i = k(x_i,x_0).
\end{equation}
Define $C_1,\ldots,C_p$ to be the ``clusters'' of the training set which satisfies that for any $i,j\in[p]$ and $a\in C_i,b\in C_j$, $i=j\iff k(x_a,x_b)=1$.
We use $\allone_S$ to denote the indicator vector, that is, $\allone_S(x) = \ind[x\in S]$.
Therefore, by definition, $K = \sum_{i=1}^p \allone_{C_i}\allone_{C_i}^\top$.
From the Sherman–Morrison formula, we can observe that the inverse of $K+n\lambda I$ should have the following form,
$$
  \tp{n\lambda I + \sum_{i=1}^p \allone_{C_i}\allone_{C_i}^\top}^{-1}
  = (n\lambda)^{-1} I + \sum_{i=1}^p \alpha_i\allone_{C_i}\allone_{C_i}^\top,
$$
where $\{\alpha_i\}_{i\in[p]}$ are coefficients that need to be determined.
Since $\inner{\allone_{C_i}}{\allone_{C_j}} = 0$ for $i\neq j$, direct calculation shows that
$$
  \tp{(n\lambda)^{-1} I + \sum_{i=1}^p \alpha_i\allone_{C_i}\allone_{C_i}^\top} \tp{n\lambda I + \sum_{i=1}^p \allone_{C_i}\allone_{C_i}^\top}
  = I + \sum_{i=1}^p ((n\lambda)^{-1}+n\alpha_i\lambda+|C_i|\alpha_i)\allone_{C_i}\allone_{C_i}^\top.
$$
Setting $\alpha_i = -(n\lambda)^{-1}(n\lambda + |C_i|)^{-1}$, we get the inverse of $K + n\lambda I$.

If $k(x_0,x_i) = 0$ for some $i\in[n]$, then by \Cref{eq:E3}, $k_{x_0} = 0$ and the prediction is $0$, as desired.

Suppose $k(x_0,x_i)=1$ for some $x_i$ and $i\in C_\ell$.
We have that by the definition of the partition, $C(x_0) = C_\ell$.
Therefore, by \Cref{eq:E3}, the prediction is as follows,
\begin{align*}
  \allone_{C_\ell}^\top \tp{(n\lambda)^{-1} I - \sum_{i=1}^p \frac{1}{n\lambda(n\lambda+|C_i|)}\allone_{C_i}\allone_{C_i}^\top}y
  &= \allone_{C_\ell}^\top \tp{(n\lambda)^{-1} I - \frac{1}{n\lambda(n\lambda+|C_\ell|)}\allone_{C_\ell}\allone_{C_\ell}^\top}y
  \\&=\tp{(n\lambda)^{-1} - \frac{|C_\ell|}{n\lambda(n\lambda+|C_\ell|)}}\allone_{C_\ell}^\top y
  \\&=\frac{\sum_{i\in C(x_0)}y_i}{n\lambda+|C(x_0)|} .
\end{align*}

\section*{Exercise 4}
1. Conditioned on $h^\star$, we have that
\begin{equation}\label{eq:E4_1}
  \Pr{\hat L_S(h) = 0} = \Pr{\forall i\in [n], h(x_i)= h^\star(x_i)}
  = (1-\Pr[x\sim \+D]{h(x)\neq h^\star(x)})^n 
  \leq \exp(-nL_{h^\star}(h)).
\end{equation}
2. Fix a hypothesis $h$ and $h^\star$.
If $L_{h^\star}(h) \leq \frac{\ln(1/\pi(h))+\ln(1/\delta)}n$, the event in the problem description holds with probability $1$.
Otherwise  by \Cref{eq:E4_1}, the following holds,
\begin{align}
  \Pr{\hat L_S(h) = 0\mid h^\star} 
  \leq\exp(-nL_{h^\star}(h))  
  \leq\exp(-(\ln(1/\pi(h))+\ln(1/\delta))) = \pi(h)\delta. \label{eq:E4_2}
\end{align}
Taking the complementary event completes the proof.
\\3. Applying the union bound over all $h$ in~\Cref{eq:E4_2}, we have that
$$
  \Pr{\exists h, \hat L_S(h) = 0 \text{ and }L_{h^\star}(h) > \frac{\ln(1/\pi(h))+\ln(1/\delta)}n \mid h^\star} \leq \delta
$$
Note that for the posterior distribution, we have that
$\Pr{h = h^\star \mid S} \propto \pi(h) \ind[\hat L_S(h) = 0]$, thus, $\hat L_S(\hat h_{\MAP}(S)) = 0$.
Therefore,
\begin{align}
  &\Pr{L_{h^\star}\tp{\hat h_{\MAP}(S)} \leq \frac{\ln(1/\pi(\hat h_{\MAP}(S)))+\ln(1/\delta)}n\mid h^\star} \nonumber
  \\\geq &\Pr{\forall h, \hat L_S(h) > 0 \text{ or }L_{h^\star}(h) \leq \frac{\ln(1/\pi(h))+\ln(1/\delta)}n \mid h^\star} \geq 1-\delta. \label{eq:E4_3}
\end{align}
4. By the maximality of $\hat h_{\MAP}(S)$, $\pi(\hat h_{\MAP}(S)) \geq \pi(h^\star)$ holds.
By \Cref{eq:E4_1}, we conclude that
\begin{align} 
  \Pr{L_{h^\star}\tp{\hat h_{\MAP}(S)} \geq \eps\mid h^\star}
  &\leq \Pr{\exists h, \hat L_S(h)=0\text{ and }L_{h^\star}\tp{h}\geq \eps\text{ and }\pi(h)\geq \pi(h^\star)\mid h^\star} \nonumber
  \\&\leq \sum_{\substack{L_{h^\star}\tp{h}\geq \eps\\\pi(h)\geq \pi(h^\star)}}\Pr{\hat L_S(h)=0\mid h^\star} \nonumber
  \\&\leq e^{-n\eps}\sum_{h:\pi(h)\geq \pi(h^\star)} \frac{\pi(h)}{\pi(h)}\leq e^{-n\eps}\sum_{h:\pi(h)\geq \pi(h^\star)} \frac{\pi(h)}{\pi(h^\star)} \leq \frac{e^{-n\eps}}{\pi(h^\star)}, \label{eq:E4_4}
\end{align}
where the second inequality follows from the union bound.
\\5. Fix $T = \ln(1/\pi(h^\star))/n$.
We have that
\begin{align} 
  \E[S]{L_{h^\star}\tp{\hat h_{\MAP}(S)}\mid h^\star} &= \int_{0}^{\infty} \Pr{L_{h^\star}\tp{\hat h_{\MAP}(S)} > t\mid h^\star} \dif t
  \leq T + \int_{T}^{\infty} \frac{e^{-nt}}{\pi(h^\star)} \dif t \nonumber
  \\&= \frac{1+\ln(1/\pi(h^\star))}{n} .\label{eq:E5_5}
\end{align}
6. By \Cref{eq:E5_5}, we conclude that
\begin{align*}
  \Pr[S,h^\star]{\hat h_{\MAP}(S)(x_0)\neq y_0} &= \E[S,h^\star,x_0]{\ind\left[\hat h_{\MAP}(S)(x_0)\neq h^\star(x_0)\right]}
  = \E[S,h^\star]{L_{h^\star}\tp{\hat h_{\MAP}(S)}}
  \\&=\E[S,h^\star]{\E{L_{h^\star}\tp{\hat h_{\MAP}(S)}\mid h^\star}} \leq \frac {1+\E[h^\star]{\ln(1/\pi(h^\star))}}{n}
  \\&= \frac{1+\-{Entropy}[\pi]}{n} .
\end{align*}
\section*{Exercise 5}
1. Direct calculation gives 
$$
  \Pr{\theta^\star = \theta\mid \{(x_i,y_i)\}_{i=1}^n} \propto \exp\tp{-\frac{\norm{\theta}_2^2}{2\lambda}}\times \exp\tp{-\sum_{i=1}^n\frac{ (\theta \cdot x_i - y_i)^2}{2\sigma^2}}.
$$
Therefore, $E(\theta) = \frac{\norm{\theta}_2^2}{2\lambda}+\sum_{i=1}^n(\theta \cdot x_i - y_i)^2/(2\sigma^2)$.
\\2. We write $X = (x_1,\ldots,x_n)$, $Y = (y_1,\ldots,y_n)^\top$ and $\hat \Sigma = \frac 1n \sum_{i=1}^n x_ix_i^\top$. 
The derivatives of $E$ are as follows,
$$
  \nabla E(\theta) = \frac{\theta}{\lambda} + \frac{XX^\top \theta - XY}{\sigma^2}\qquad\text{and}\qquad
  \nabla^2 E(\theta) = \frac{I_d}{\lambda} + \frac n{\sigma^2}\hat \Sigma.
$$
Since $\hat \theta_\MAP$ minimizes $E$, $\nabla E(\theta) = 0$ implies
$$
  \hat \theta_\MAP = \tp{\frac {\sigma^2 I_d}{\lambda} + XX^\top}^{-1}XY .
$$ 
Therefore, the Laplace approximation is as follows,
$$
  \Pr{\theta^\star = \theta\mid \{(x_i,y_i)\}_{i=1}^n} \propto \exp\tp{-\frac12(\theta - \hat \theta_\MAP)^\top\tp{\frac{I_d}{\lambda} + \frac{n}{\sigma^2}\hat\Sigma}(\theta-\hat \theta_\MAP)} .
$$
In this setting, the mean and covariance are as follows,
$$
  \mu = \tp{\frac {\sigma^2 I_d}{\lambda} + XX^\top}^{-1}XY,\qquad \Sigma = \tp{\frac{I_d}{\lambda} + \frac n{\sigma^2}\hat \Sigma}^{-1} .
$$
3. Direct calculation gives 
$$
  \Pr{\theta^\star = \theta\mid \{(x_i,y_i)\}_{i=1}^n} \propto \exp\tp{-\frac{\norm{\theta}_2^2}{2\lambda}}\times \prod_{i=1}^n\sigma\tp{ y_i\inner{\theta}{x_i}} .
$$
Therefore, $E(\theta) = \frac{\norm{\theta}_2^2}{2\lambda}-\sum_{i=1}^n \ln(\sigma\tp{ y_i\inner{\theta}{x_i}})$.
\\4. The derivatives of $E$ are as follows,
$$
  \nabla E(\theta) = \frac{\theta}{\lambda} - \sum_{i=1}^n y_i \sigma\tp{-y_i\inner{x_i}{\theta}}x_i
  \qquad\text{and}\qquad 
  \nabla^2 E(\theta) = \frac{I_d}{\lambda} + \sum_{i=1}^n  \sigma\tp{\inner{x_i}{\theta}}\sigma\tp{-\inner{x_i}{\theta}}x_ix_i^\top .
$$
The maximizer $\hat \theta_\MAP$ satisfies $\nabla E(\hat \theta_\MAP) = 0$.
For the Laplace approximation, we have that the posterior distribution is proportional to 
$$
  \exp\tp{-\frac12(\theta - \hat \theta_\MAP)^\top\tp{\frac{I_d}{\lambda} + \sum_{i=1}^n  \sigma\tp{\inner{x_i}{\hat \theta_\MAP}}\sigma\tp{-\inner{x_i}{\hat \theta_\MAP}}x_ix_i^\top}(\theta-\hat \theta_\MAP)} .
$$
In this setting, the mean and covariance are as follows,
$$
  \mu = \hat \theta_\MAP,\qquad \Sigma = \tp{\frac{I_d}{\lambda} + \sum_{i=1}^n  \sigma\tp{\inner{x_i}{\hat \theta_\MAP}}\sigma\tp{-\inner{x_i}{\hat \theta_\MAP}}x_ix_i^\top}^{-1} .
$$
\end{document}